\documentclass[11pt,a4paper]{article}

\usepackage[utf8]{inputenc}
\usepackage[T1]{fontenc}
\usepackage{lmodern}
\usepackage{microtype}
\usepackage{geometry}
\usepackage{amsmath}
\usepackage{amssymb}
\usepackage{xcolor}
\usepackage{listings}
\usepackage{enumitem}
\usepackage[hidelinks]{hyperref}

\geometry{margin=1in}
\setlist{nosep}
\setlength{\parskip}{0.55em}
\setlength{\parindent}{0pt}

\definecolor{codebackground}{RGB}{248,248,248}
\definecolor{codecomment}{RGB}{90,110,90}
\definecolor{codekeyword}{RGB}{20,70,160}
\definecolor{codestring}{RGB}{150,60,40}

\lstdefinestyle{pythonstyle}{
  language=Python,
  basicstyle=\ttfamily\small,
  keywordstyle=\color{codekeyword}\bfseries,
  commentstyle=\color{codecomment},
  stringstyle=\color{codestring},
  backgroundcolor=\color{codebackground},
  frame=single,
  rulecolor=\color{black!15},
  breaklines=true,
  breakatwhitespace=true,
  showstringspaces=false,
  columns=fullflexible,
  keepspaces=true,
  tabsize=4
}

\lstdefinestyle{textstyle}{
  basicstyle=\ttfamily\small,
  backgroundcolor=\color{codebackground},
  frame=single,
  rulecolor=\color{black!15},
  breaklines=true,
  columns=fullflexible,
  keepspaces=true
}

\title{From Borges to Meaning: Progressive Reduction of the Library of Babel}
\author{}
\date{}

\begin{document}

\maketitle

\tableofcontents
\newpage

\phantomsection
\section*{Abstract}
\addcontentsline{toc}{section}{Abstract}

This article investigates a sequence relating to mathematical transformations of \emph{The Library of Babel} originally proposed by Jorge Luis Borges. Starting from the classical combinatoric construction containing every possible sequence of characters, we progressively constrain the Library so that:

\begin{enumerate}
  \item the number of possible books decreases;
  \item the average semantic density increases;
  \item random pages become increasingly human-readable.
\end{enumerate}

At each stage we introduce:

\begin{itemize}
  \item a mathematical model;
  \item combinatoric analysis;
  \item asymptotic reduction;
  \item Python algorithms that are capable of generating pages on demand from compact seeds.
\end{itemize}

The experiment reveals a transition from pure entropy toward structured language and eventually toward semantic manifolds approximating meaningful literature.

\section{Original Borges Library}

Borges defines books with:

\begin{itemize}
  \item 410 pages;
  \item 40 lines per page;
  \item 80 characters per line;
  \item an alphabet of 25 symbols.
\end{itemize}

Total character positions:

\[
C = 410 \times 40 \times 80
\]

\[
410 \times 40 \times 80 = 1{,}312{,}000
\]

Total number of books:

\[
25^{1{,}312{,}000}
\]

Approximation:

\[
25^{1{,}312{,}000}
\approx
1.95 \times 10^{1{,}834{,}097}
\]

This construction maximizes entropy but minimizes meaning.

Almost every book is unreadable noise.

\section{Word-Based Library}

\subsection{Motivation}

Human language is not constructed from arbitrary character sequences.

Replacing characters with words immediately injects semantic structure.

Assume:

\begin{itemize}
  \item vocabulary size:
  \[
  W = 100{,}000
  \]
  \item punctuation symbols:
  \[
  P = 6
  \]
  \item average token width $\approx 6$ characters.
\end{itemize}

Approximate token slots:

\[
N \approx \frac{1{,}312{,}000}{6}
\]

\[
N \approx \frac{1{,}312{,}000}{6} \approx 218{,}667
\]

Total books:

\[
(100{,}006)^{218{,}667}
\]

Approximation:

\[
\approx 10^{1{,}093{,}340}
\]

The Library shrinks enormously while readability increases dramatically.

\subsection{Python Generator}

\begin{lstlisting}[style=pythonstyle]
import hashlib

VOCABULARY = ["the", "house", "river", "memory", "darkness"]

def token_at(seed: str, position: int):
    data = f"{seed}:{position}".encode("utf-8")
    digest = hashlib.sha256(data).digest()
    value = int.from_bytes(digest[:8], "big")
    return VOCABULARY[value % len(VOCABULARY)]

def generate_page(seed, page, tokens_per_page=300):
    start = page * tokens_per_page
    return [
        token_at(seed, i)
        for i in range(start, start + tokens_per_page)
    ]
\end{lstlisting}

The complete book never needs to exist in memory.

\section{Punctuation-Constrained Library}

\subsection{Constraint}

No two punctuation marks may appear consecutively.

This removes sequences like:

\begin{lstlisting}[style=textstyle]
word . ? , !
\end{lstlisting}

which are structurally invalid.

\subsection{Combinatorics}

If exactly $k$ punctuation marks appear:

\[
\binom{N-k+1}{k}
\]

ways exist to place them without adjacency.

Total Library:

\[
\sum_{k=0}^{\lfloor (N+1)/2 \rfloor}
\binom{N-k+1}{k} P^k W^{N-k}
\]

This slightly reduces entropy while improving grammatical plausibility.

\subsection{Python Implementation}

\begin{lstlisting}[style=pythonstyle]
def valid_sequence(tokens):
    punct = {".", ",", ";", ":", "?", "!"}

    for i in range(len(tokens) - 1):
        if tokens[i] in punct and tokens[i+1] in punct:
            return False

    return True
\end{lstlisting}

\section{Sentence-Structured Library}

\subsection{Constraint}

Every sentence must contain:

\begin{itemize}
  \item exactly 15 words;
  \item followed by one punctuation mark.
\end{itemize}

Pattern:

\begin{lstlisting}[style=textstyle]
W W W W W W W W W W W W W W W P
\end{lstlisting}

Sentence count:

\[
S = \left\lfloor \frac{N}{16} \right\rfloor
\]

Total Library:

\[
W^{15S} P^S
\]

Approximation:

\[
\approx 10^{1{,}035{,}639}
\]

The Library loses vast amounts of entropy while gaining strong textual regularity.

Random pages now resemble primitive prose.

\subsection{Python Generator}

\begin{lstlisting}[style=pythonstyle]
WORDS_PER_SENTENCE = 15

PUNCT = [".", "?", "!"]

def sentence(seed, index):
    words = [
        token_at(seed, index * 16 + i)
        for i in range(WORDS_PER_SENTENCE)
    ]

    punct = PUNCT[
        hash(f"{seed}:{index}") % len(PUNCT)
    ]

    return " ".join(words) + punct
\end{lstlisting}

\section{Grammar-Constrained Library}

\subsection{Constraint}

Sentences must satisfy grammatical templates.

Example:

\begin{lstlisting}[style=textstyle]
DET ADJ NOUN VERB DET NOUN .
\end{lstlisting}

Define vocabulary partitions:

\begin{itemize}
  \item nouns;
  \item verbs;
  \item adjectives;
  \item adverbs;
  \item determiners;
  \item pronouns;
  \item prepositions.
\end{itemize}

\subsection{Combinatorics}

If:

\[
D, A, N, V
\]

represent category counts, then one sentence template produces:

\[
D \cdot A \cdot N \cdot V \cdot D \cdot N \cdot P
\]

possible sentences.

A book with $S$ such sentences gives:

\[
(DANVDNP)^S
\]

This reduction is enormous.

However, readability increases dramatically.

Random pages now resemble machine-generated language.

\subsection{Python Implementation}

\begin{lstlisting}[style=pythonstyle]
GRAMMAR = [
    ("DET", ["the", "a"]),
    ("ADJ", ["dark", "ancient"]),
    ("NOUN", ["river", "library", "mirror"]),
    ("VERB", ["contains", "reflects"]),
    ("DET", ["the", "a"]),
    ("NOUN", ["truth", "memory"]),
]

def grammar_sentence(seed, sentence_id):
    words = []

    for i, (_, vocab) in enumerate(GRAMMAR):
        token = token_at(seed, sentence_id * 100 + i)
        words.append(vocab[hash(token) % len(vocab)])

    return " ".join(words) + "."
\end{lstlisting}

\section{Semantic-Constrained Library}

\subsection{Constraint}

Not all grammatically valid sentences are meaningful.

Introduce semantic neighborhoods.

Example:

\begin{lstlisting}[style=textstyle]
river <-> water <-> flow <-> current
library <-> books <-> shelves <-> archive
\end{lstlisting}

Words may only follow semantically compatible words.

This transforms the Library from a combinatoric explosion into a constrained semantic graph.

\subsection{Markov Formulation}

Let:

\[
T_{ij}
\]

be transition probability from token $i$ to token $j$.

Then valid books are paths through the semantic graph.

Approximate sequence count:

\[
\lambda_{\max}^{N}
\]

where:

\[
\lambda_{\max}
\]

is the dominant eigenvalue of the transition matrix.

This is a profound reduction.

Meaning emerges from spectral constraints.

\subsection{Python Implementation}

\begin{lstlisting}[style=pythonstyle]
import random

GRAPH = {
    "river": ["water", "flow", "stone"],
    "water": ["current", "river"],
    "library": ["books", "archive"],
}

def semantic_walk(seed, start, length):
    random.seed(seed)

    word = start
    result = [word]

    for _ in range(length - 1):
        next_words = GRAPH.get(word, [word])
        word = random.choice(next_words)
        result.append(word)

    return result
\end{lstlisting}

\section{Topic-Coherent Library}

\subsection{Constraint}

Entire books must remain inside a thematic manifold.

Examples:

\begin{itemize}
  \item cosmology;
  \item grief;
  \item mathematics;
  \item theology;
  \item medieval warfare.
\end{itemize}

This introduces long-range coherence.

Now the Library approximates actual literature.

\subsection{Information-Theoretic Interpretation}

The original Borges Library maximizes entropy:

\[
H_{\max}
\]

Every subsequent constraint reduces entropy:

\[
H_1 > H_2 > H_3 > H_4
\]

while increasing mutual information:

\[
I(\text{text};\text{meaning})
\]

The experiment demonstrates that human-readable language occupies an infinitesimally thin manifold inside the full combinatoric space of possible books.

\section{Final Observation}

The original Library is not terrifying because it contains all books.

It is terrifying because meaningful books are drowned inside an ocean of entropy.

Each constraint that is introduced in this article:

\begin{itemize}
  \item decreases combinatoric volume;
  \item increases semantic density;
  \item and moves the Library closer to the tiny structured region inhabited by human thought.
\end{itemize}

The true Library of meaning is not infinite.

It is an extraordinarily compressed island hidden inside Borges' cosmic desert of symbols.

\end{document}
